\section{Data}
\label{sec:data}

\subsection{Source and Inventory}

Describe where the data comes from, how much there is, and its properties
(resolution, duration, dates). \Cref{tab:inventory} summarises the data.

% ---- Table: caption ABOVE the table ----
\begin{table}
  \caption{Short title of the table. Add one sentence of context if needed.}
  \label{tab:inventory}
  \begin{autotable}{L R R}
    \toprule
    \textbf{Item} & \textbf{Count} & \textbf{Subset} \\
    \midrule
    Item A & 100 & 20 \\ \rowline
    Item B & 50  & 5  \\ \rowline
    Item C & 10  & 0  \\
    \bottomrule
  \end{autotable}
\end{table}

Describe what makes the data difficult, with examples in
\cref{fig:example-a,fig:example-b,fig:example-c,fig:example-d}.

% ---- Sub-figures: "Figure 1a", "Figure 1b", ... caption BELOW each ----
\begin{figure}
  \begin{subfigure}{\linewidth}
    \includegraphics[width=\linewidth, height=4.5cm]{example-image-a}
    \caption{First example. One sentence on what it shows.}
    \label{fig:example-a}
  \end{subfigure}

  \subfiguregap
  \begin{subfigure}{\linewidth}
    \includegraphics[width=\linewidth, height=4.5cm]{example-image-b}
    \caption{Second example.}
    \label{fig:example-b}
  \end{subfigure}
\end{figure}

% Same figure continued on the next page (keeps the number, goes on to c, d)
\begin{figure}
  \ContinuedFloat
  \begin{subfigure}{\linewidth}
    \includegraphics[width=\linewidth, height=4.5cm]{example-image-c}
    \caption{Third example.}
    \label{fig:example-c}
  \end{subfigure}

  \subfiguregap
  \begin{subfigure}{\linewidth}
    \includegraphics[width=\linewidth, height=4.5cm]{example-image}
    \caption{Fourth example.}
    \label{fig:example-d}
  \end{subfigure}
\end{figure}

\subsection{Train/Val/Test Split}

Explain how the data is split and why the split avoids leakage between
training and evaluation (\cref{tab:split}).

\begin{table}
  \caption{Split summary.}
  \label{tab:split}
  \begin{autotable}{L R R R}
    \toprule
    \textbf{Split} & \textbf{Groups} & \textbf{Items} & \textbf{Frames} \\
    \midrule
    Train & 0 & 0 & 0 \\ \rowline
    Val   & 0 & 0 & 0 \\ \rowline
    Test  & 0 & 0 & 0 \\
    \bottomrule
  \end{autotable}
\end{table}

\subsection{Labels}

Describe how labels were made and checked.

\subsubsection{Labeling Rules}

Use a third heading level only when a subsection needs its own parts.
